\section{Methodology}
\label{sec:methodology}

\subsection{Vendor-Portable Code Path}
\label{sec:methodology:portability}

All trainer scripts use \texttt{trainer/device\_utils.py} as the single point of
contact for hardware-specific behavior.  The module exposes two layers:

\textbf{Free functions} answer environment-level questions:
\texttt{is\_rocm()}, \texttt{is\_cuda\_native()}, \texttt{vendor\_name()},
\texttt{dist\_backend()}, \texttt{default\_device\_str()},
\texttt{set\_current\_device()}, and \texttt{print\_device\_info()}.
\texttt{dist\_backend()} returns \texttt{"nccl"} on both CUDA and ROCm because
PyTorch-ROCm exposes RCCL under that string alias.
\texttt{print\_device\_info()} logs a one-line banner at startup
(Figure~\ref{fig:banner}) that is captured for every experiment run.

\textbf{\texttt{DeviceCtx}} (frozen dataclass) encapsulates per-trainer state.
Each trainer builds it once:
\begin{lstlisting}[language=Python,basicstyle=\small\ttfamily]
dev = DeviceCtx.from_arg(args.device)
if dist.is_initialized():
    dev = dev.for_rank(local_rank)
autocast_ctx = dev.autocast(args.dtype)   # nullcontext on CPU
scaler       = dev.grad_scaler(enabled=(args.dtype == 'float16'))
\end{lstlisting}
\texttt{dev.autocast()} dispatches to
\texttt{torch.amp.autocast(device\_type="cuda", ...)} for both NVIDIA and AMD
(ROCm aliases \texttt{cuda} in its PyTorch build); it returns
\texttt{nullcontext()} on CPU so no special-casing is needed.

\textbf{Consequence for this study.}  The only files that differ between a CUDA
run and a ROCm run are the PyTorch wheel and the \texttt{HSA\_OVERRIDE\_GFX\_VERSION}
environment variable (required for gfx1100-family targets; not needed for MI350
gfx940).  Trainer logic, model code, checkpoint format, and dataset loading are
bit-for-bit identical.

\subsection{Reproducibility Recipe}
\label{sec:methodology:repro}

\begin{itemize}[leftmargin=*]
  \item \textbf{Software pins.}  All experiments use:
    \texttt{torch==2.10.0+cu12x} (NVIDIA) and
    \texttt{torch==2.10.0+rocm7.1} (AMD);
    \texttt{transformers==4.5x} (pinned; see Appendix~\ref{sec:appendix:install});
    \texttt{vllm} and \texttt{sglang} at the commit SHAs listed in
    Table~\ref{tab:versions}.
  \item \textbf{Random seeds.}  Every trainer calls
    \texttt{setup\_seed(42 + rank)} at startup, which seeds Python
    \texttt{random}, NumPy, PyTorch, and \texttt{torch.cuda} (CUDA and ROCm
    use the same API surface).
  \item \textbf{Data.}  Datasets are downloaded at fixed versions; SHA-256
    hashes are recorded in Appendix~\ref{sec:appendix:data}.
  \item \textbf{Banner capture.}  Every run is launched via a wrapper that
    captures the first \texttt{[device] ...} line from
    \texttt{print\_device\_info()} and saves it alongside metrics.  This guards
    against mis-scheduled jobs.
  \item \textbf{Checkpoint naming.}  MiniMind uses load-bearing checkpoint
    names of the form
    \texttt{\{save\_weight\}\_\{hidden\_size\}\{\_moe?\}.pth}.
    The same naming convention is enforced in all experiment scripts.
\end{itemize}

\subsection{Model-FLOP Utilization}
\label{sec:methodology:mfu}

Following~\citet{chowdhery2023palm}, we define Model-FLOP Utilization (MFU) as:
\[
  \text{MFU} = \frac{\text{achieved FLOPs/s}}{\text{peak FLOPs/s of the device}}
\]
where \emph{achieved FLOPs/s} $= C \times S / T$,
with $C$ the FLOPs per token, $S$ the number of tokens processed in interval
$T$.  For a MiniMind forward + backward pass on a dense model:
\[
  C_{\text{dense}} = 6 \times P
\]
where $P$ is the number of trainable parameters (the factor of 6 accounts for
one forward and two backward passes, each costing $\approx 2P$
multiply-accumulates per token~\citep{chowdhery2023palm}).

For the MoE variant, only the active experts contribute to the compute per
token, so:
\[
  C_{\text{MoE}} = 6 \times P_{\text{active}}
\]
where $P_{\text{active}} = P_{\text{base}} + n_{\text{active}} \times
P_{\text{expert}}$, with $n_{\text{active}}$ the number of experts activated
per token (default 1 in MiniMind).

Peak FLOPs/s: H100 SXM5 delivers 989\,TFLOPs/s (bf16 tensor core);
MI350 delivers \PH{peak TFLOPs/s from AMD spec sheet}.
MFU is computed offline from logged step times by
\texttt{paper/experiments/compute\_mfu.py}.

\subsection{Inference Metrics}
\label{sec:methodology:inf}

\begin{itemize}[leftmargin=*]
  \item \textbf{TTFT} (time-to-first-token): wall time from request arrival to
    first output token, measured at the server level for
    \texttt{serve\_openai\_api.py} and at the engine level for vLLM / SGLang.
  \item \textbf{Inter-token latency} (ITL): wall time between consecutive output
    tokens during decode; reported as p50 and p95 over a 200-request run.
  \item \textbf{Tokens/s}: aggregate decode throughput across concurrent
    requests.
  \item \textbf{Requests/s}: completed requests per second at each concurrency
    level.
  \item \textbf{HBM at steady state}: \texttt{torch.cuda.max\_memory\_allocated()}
    sampled once per second during the benchmark window.
  \item \textbf{Prefill/decode split}: derived by logging TTFT separately from
    total generation time.
\end{itemize}

\subsection{Experimental Configurations}
\label{sec:methodology:configs}

Table~\ref{tab:configs} summarizes the experimental configurations used
throughout this paper.  For training, we use the ``mini'' datasets shipped with
the MiniMind repository, which are sufficient to reproduce the upstream
single-3090 reference numbers and yield comparable relative performance across
platforms.  Full dataset hashes are in Appendix~\ref{sec:appendix:data}.

\begin{table}[h]
\centering
\caption{Summary of experimental configurations.}
\label{tab:configs}
\small
\begin{tabular}{llll}
\toprule
\textbf{Dimension} & \textbf{Values} & \textbf{Primary} & \textbf{Note} \\
\midrule
Model variant & Dense, MoE & Dense & hidden\_size=768, 8 layers \\
\# GPUs & 1, 4, 8 & 8 & Single node, DDP \\
dtype & bf16, fp16 & bf16 & fp16 uses GradScaler \\
\texttt{torch.compile} & 0, 1 & 0 & LoRA: auto-disabled \\
Rollout engine & torch, sglang & torch & PPO/GRPO only \\
Inference engine & HF, vLLM, SGLang & vLLM & Block E \\
\bottomrule
\end{tabular}
\end{table}
